\section{A polynomial lower bound for real ray-class values}
\label{sec:ray-real-gap}

The real exceptional zero in
Section~\ref{sec:level-prime-selection} requires a polynomial separation
from one. We prove the quantitative implication from the ray-family
counting and theta identities. All constants below are independent of
the class group and its modulus.

Let \(G\) be a finite abelian group of order \(H\), let \(D\geq4\),
\(H\leq D\), and let \(a_g(n)\geq0\) satisfy
\[
 \begin{gathered}
 a_g(0)=0,\qquad a_0(1)=1,\qquad a_g(1)=0\quad(g\ne0),\\
 D^{-1}\leq\kappa\leq4,\qquad
 \left|\sum_{n\leq u}a_g(n)-\kappa u\right|
                                  \leq10\sqrt u\quad(u\geq1),\\
 \left|\sum_{n\geq1}a_g(n)e^{-tn}-\frac{\kappa}{t}\right|
                                   \leq D^{10}t\quad(0<t\leq1).
 \end{gathered}
\]
For every real character \(\chi:G\to\{1,-1\}\), put
\(a_\chi(n)=\sum_g\chi(g)a_g(n)\).
We require the positivity furnished by the ray-family Euler product,
in its precise coefficient form
\begin{equation}
 \label{eq:ray-convolution-positive}
          \sum_{mn=k}a_1(m)a_\chi(n)\geq0\qquad(k\geq1).
\end{equation}
The subscript \(1\) denotes the trivial character. These are exactly
the counting, theta, and real-convolution properties used from
\srcid{LFam}. This section proves their analytic consequence. Section~\ref{sec:quadratic-ray-family} constructs them from the ideals
of the quadratic field.

\subsection{Continuation and uniform bounds near one}

Nonnegativity and the counting bound give, for \(n\geq1\),
\[
       a_g(n)\leq\kappa n+10\sqrt n\leq14n,
                         \qquad |a_\chi(n)|\leq14Dn.
\]
Let \(\rho=\kappa H\), and write
\[
 A_\chi(u)=\sum_{n\leq u}a_\chi(n),\qquad
 R_\chi(u)=A_\chi(u)-{\bf1}_{\chi=1}\rho u .
\]
Orthogonality gives \(|R_\chi(u)|\leq10D\sqrt u\).
Consequently
\[
        I_\chi(z)=\int_1^\infty R_\chi(u)u^{-z-1}\,du
\]
is holomorphic for \(\Re z>1/2\), with
\[
 |I_\chi(z)|\leq\frac{10D}{\Re z-1/2},\qquad
 |I'_\chi(z)|\leq\frac{10D}{(\Re z-1/2)^2}.
\]
Both assertions follow by domination on every closed subregion of
that half-plane. Differentiation inserts \(-\log u\), and
\(\int_1^\infty u^{-a-1}\log u\,du=a^{-2}\) for \(a>0\).

For \(\Re z>1\), summation by parts gives
\[
       L_\chi(z)
          ={\bf1}_{\chi=1}\frac{\rho z}{z-1}
                              +zI_\chi(z).
\]
Thus the regularized function used in the source is
\[
 F_\chi(z)=
 \begin{cases}
   \rho z+(z-1)zI_\chi(z),&\chi=1,\\
   zI_\chi(z),&\chi\ne1.
 \end{cases}
\]
It is holomorphic for \(\Re z>1/2\), and \(F_1(1)=\rho\).

For \(|z-1|\leq1/4\), we have
\(\Re z\geq3/4\), \(|z|\leq5/4\),
\(|I_\chi(z)|\leq40D\), and \(|I'_\chi(z)|\leq160D\).
If \(\chi\ne1\), differentiation gives
\[
                   |F'_\chi(z)|\leq40D+(5/4)160D=240D.
\]
For the trivial character, use
\[
 F'_1(z)=\rho+(2z-1)I_1(z)+(z-1)zI'_1(z).
\]
Since \(\rho\leq4D\), its norm is at most
\(4D+(3/2)40D+(5/16)160D=114D\).
In particular the weaker uniform bound used by the formal argument holds:
\begin{equation}
 \label{eq:ray-derivative-bound}
                 |F'_\chi(z)|\leq18000D
                             \qquad(|z-1|\leq1/4).
\end{equation}

\subsection{The smoothed positive kernel}

For \(y>0\) define
\[
            W(y)=\int_0^\infty e^{-y/w-w}\,\frac{dw}{w}.
\]
It is positive and decreasing in \(y\).
The elementary inequality
\(e^{-y/w}\leq6w^3/y^3\) gives
\[
           0<W(y)\leq\frac6{y^3}
                       \int_0^\infty w^2e^{-w}\,dw
                     =\frac{12}{y^3}.
\]
In particular \(w_0=W(1)>0\) is an absolute finite constant.

Take a nontrivial real character \(\chi\), and set
\[
       P(t)=\sum_{n\geq1}a_1(n)e^{-tn},\qquad
       Q(t)=\sum_{n\geq1}a_\chi(n)e^{-tn}.
\]
For \(0<t\leq1\), summing the theta bounds and using character
orthogonality gives
\[
       |P(t)-\rho/t|\leq D^{11}t,\qquad
                          |Q(t)|\leq D^{11}t.
\]
For \(t\geq1\), the coefficient bound gives
\[
 |P(t)|,\ |Q(t)|
       \leq14D\sum_{n\geq1}ne^{-nt}
       =14D\frac{e^{-t}}{(1-e^{-t})^2}.
\]
There is therefore an absolute constant \(C\geq1\), for instance
\(C=100\), for which, with \(E=CD^{11}\),
\begin{equation}
 \label{eq:ray-theta-kernel-bounds}
 \begin{split}
       |P(t)-\rho/t|&\leq E\min(t,t^{-1}),\\
                    |Q(t)|&\leq E\min(t,t^{-1}),\\
       |Q(t)|&\leq CD e^{-t/2}\quad(t\geq1).
 \end{split}
\end{equation}
For the first two inequalities on \(t\geq1\), multiply the
preceding exponential estimate by \(t\) and use \(te^{-t}\leq1\),
together with \(\rho\leq4D\).
These bounds prove that \(Q\) is integrable at zero and infinity.

For \(v>0\), expansion under the integral gives
\begin{equation}
 \label{eq:ray-kernel-product}
 J(v):=\int_0^\infty P(v/w)Q(w)\,\frac{dw}{w}
           =\sum_{m,n\geq1}a_1(m)a_\chi(n)W(vmn).
\end{equation}
Indeed the substitution \(u=nw\) transforms
\(\int e^{-vm/w-nw}\,dw/w\) to \(W(vmn)\).
Every exchange is justified absolutely by
\[
 \sum_{m,n\geq1}|a_1(m)a_\chi(n)|W(vmn)
       \leq\frac{12(14D)^2}{v^3}
                              \sum_{m,n\geq1}\frac1{m^2n^2}<\infty .
\]
Grouping by \(mn=k\) and applying
\eqref{eq:ray-convolution-positive} yields
\[
                         J(v)\geq W(v),
\]
because the coefficient at \(k=1\) is one.

Let
\[
       \omega(w)=\frac{\min(w,w^{-1})}{w\sqrt w},\qquad
       K=\int_0^\infty\omega(w)\,dw=\frac83 .
\]
Write \(P(v/w)=\rho w/v+R_v(w)\).
The inequality \(\min(u,u^{-1})\leq\sqrt u\), for \(u>0\),
and \eqref{eq:ray-theta-kernel-bounds} give
\[
       \left|\frac{R_v(w)Q(w)}w\right|
                         \leq E^2\sqrt v\,\omega(w).
\]
The right side is integrable. Thus, putting
\(L=\int_0^\infty Q(w)\,dw\), we obtain
\begin{equation}
 \label{eq:ray-upper-kernel}
                  J(v)\leq\frac{\rho L}{v}+E^2\sqrt v K .
\end{equation}
This also gives an independent integrable majorant for the
decomposition of the integral.

Set
\[
 \delta=\min\left(1,\frac{w_0}{2(C^2K+1)}\right),\qquad
                         s=\delta D^{-22},\qquad v=s^2 .
\]
Then \(0<v\leq1\), so \(W(v)\geq w_0\), and
\[
            E^2\sqrt v K=C^2D^{22}sK=C^2\delta K\leq w_0/2 .
\]
Combining the lower and upper estimates for \(J(v)\) gives
\[
                \frac{w_0}{2}\leq\frac{\rho L}{s^2},
                    \qquad
                L\geq\frac{w_0\delta^2}{8D^{45}},
\]
where the second inequality uses \(0<\rho\leq4D\).

\subsection{Identifying the value at one}

For \(\Re z>2\), absolute convergence and the gamma integral give
\[
       \int_0^\infty Q(t)t^{z-1}\,dt
                   =\Gamma(z)L_\chi(z).
\]
The coefficient bound \(14Dn\) justifies the exchange of sum and
integral in that region. The small-\(t\) bound \(Q(t)=O(D^{11}t)\)
and the exponential bound at infinity show that the left side is
holomorphic for \(\Re z>-1\), by domination on compact subsets.
The gamma function is holomorphic and nonzero for \(\Re z>0\).
The identity theorem, applied in \(\Re z>1/2\), therefore identifies
this integral divided by \(\Gamma(z)\) with the continuation
\(F_\chi(z)\) there. At \(z=1\), since \(\Gamma(1)=1\), it gives
\[
                           F_\chi(1)=L.
\]
For the trivial character \(F_1(1)=\rho\geq D^{-1}\).
Consequently the positive absolute constant
\[
                  c_0=\min(1,w_0\delta^2/8)
\]
satisfies, for every real character including the trivial one,
\begin{equation}
 \label{eq:ray-real-value}
                         |F_\chi(1)|\geq c_0D^{-45}.
\end{equation}
If \(3/4\leq\beta<1\) and \(F_\chi(\beta)=0\), integrate
\(F'_\chi\) along the segment from \(\beta\) to \(1\) and use
\eqref{eq:ray-derivative-bound}. Then
\[
       c_0D^{-45}\leq|F_\chi(1)|
                         \leq18000D(1-\beta),
 \qquad
             1-\beta\geq\frac{c_0}{18000}D^{-46}.
\]
This supplies the required polynomial gap with an explicit
admissible exponent. The ray-family identities proved in Section~\ref{sec:quadratic-ray-family}
and the uniform characterwise explicit formula enter the norm-prime count. The lower bound at one and the
deduction of the real-zero gap have been proved here.

\paragraph{Source map.}
The exponential estimates correspond to
\srcid{ExpSumBounds}; the kernel and its positive convolution to
\srcid{NormRayClassFamilyFOneGeOfRealCharacter}; and the last integration
to \srcid{ExistsRayClassLFunctionRealZeroFarFromOne}.
\footnote{\href{https://github.com/Deicyde/riemann-hypothesis/blob/40e56bab026552e089d7afcfeff08d2defb9c005/src/Development/AbelianVarietyTheory/Definitions/GlobalTwisting/LamIdealOfNotPrimePow.lean}{The ray-family hypotheses};
\href{https://github.com/Deicyde/riemann-hypothesis/blob/40e56bab026552e089d7afcfeff08d2defb9c005/src/Development/AbelianVarietyTheory/Proofs/GlobalInterpolation/NormRayClassFamilyFOneGeOfRealCharacter.lean}{the positive-kernel proof};
\href{https://github.com/Deicyde/riemann-hypothesis/blob/40e56bab026552e089d7afcfeff08d2defb9c005/src/Development/AbelianVarietyTheory/Proofs/GlobalInterpolation/ExistsRayClassLFunctionRealZeroFarFromOne.lean}{the real-zero gap}.}
